% SPDX-License-Identifier: CC-BY-4.0
% English commentary accompanying LocalSIC12.tex, by Chatgpt.
\documentclass[11pt,a4paper]{article}
\usepackage[T1]{fontenc}
\usepackage{lmodern,microtype,amsmath,amssymb,booktabs}
\usepackage[top=20mm,bottom=21mm,left=25mm,right=25mm]{geometry}
\usepackage{xcolor}
\usepackage[colorlinks=true,linkcolor=blue!45!black,citecolor=blue!45!black,urlcolor=blue!45!black]{hyperref}
\hypersetup{pdftitle={Commentary on the local SIC obstruction in dimension twelve},pdfauthor={Chatgpt}}
\newcommand{\tr}{\operatorname{tr}}
\newcommand{\C}{\mathbb C}
\newcommand{\Pset}{\mathcal P}
\setlength{\parindent}{1.2em}
\begin{document}
\begin{center}
{\large\itshape Commentary}\par\smallskip
{\LARGE\bfseries The local SIC obstruction in dimension twelve\\[2pt]
and the spectra of signed Pauli sums}\par\medskip
Chatgpt\quad$\cdot$\quad 27 November 2026
\end{center}

The manuscript \emph{Abelian covariance of symmetric informationally complete measurements in dimension twelve}~\cite{sic} proves that the local displacements on two qubits and a qutrit cannot generate a symmetric informationally complete measurement (SIC). Banacki et al.\ had explicitly posed this existence question~\cite[Sec.~5.1]{banacki}. Grassl's exact solution in dimension twelve uses the ordinary twelve-level Weyl--Heisenberg action~\cite{grassl}. The local action replaces the single twelve-level shift and clock with separate displacements on each of the three factors.

A subsequent literature and code search identified a second route to the local obstruction. The \emph{StabilizerHalf} project by Wangwy-sjtu~\cite{half} computes the maximum half-order stabilizer R\'enyi entropy of a pure two-qubit state. Its exact table of characteristic polynomials implies the local SIC result once one uses the qutrit to bound the rank of a two-qubit marginal.

\section*{A forbidden eigenvalue}
Let $\Pset=\{I,X,Y,Z\}^{\otimes2}\setminus\{I\otimes I\}$, the fifteen nonidentity Hermitian Pauli operators on the two qubits. A SIC consists of $d^2$ rays with squared overlaps $1/(d+1)$ between distinct rays. In an orbit generated by local displacements from a unit vector $\psi\in\C^4\otimes\C^3$, this would require
\[
  |r_P|=1/\sqrt{13},\qquad
  r_P=\langle\psi|(P\otimes I_3)|\psi\rangle,\qquad P\in\Pset.
\]
The $r_P$ are real. Pauli orthogonality gives the expansion of the two-qubit marginal $\rho$,
\[
  \rho=\tfrac14\bigl(I+S/\sqrt{13}\bigr),\qquad
  S=\sum_{P\in\Pset}s_PP,\qquad s_P\in\{-1,1\}.
\]
The amplitude matrix of $\psi$ has four rows and three columns, so $\rho$ has rank at most three. Its kernel therefore supplies an eigenvector of $S$ with eigenvalue $-\sqrt{13}$.

For each choice of signs, the characteristic polynomial of $S$ is
$f(x)=x^4-30x^2+Bx+D$. The $x^3$ and $x^2$ coefficients follow from $\tr S=0$ and $\tr S^2=60$. The table in \texttt{check.py}~\cite{half} gives all possible integer pairs $(B,D)$. The counts in the last row apply separately to each sign of $B$.
\begin{center}
\begin{tabular}{lrrrrrr}
\toprule
$|B|$ & 72 & 56 & 40 & 24 & 8 & 8\\
$D$ & $-27$ & $-27$ & 5 & 69 & 37 & 165\\
Sign assignments for each $(B,D)$ & 640 & 1,920 & 2,304 & 3,840 & 5,760 & 1,920\\
\bottomrule
\end{tabular}
\end{center}
These twelve classes exhaust the $2^{15}=32{,}768$ assignments and agree with the six pairs $(|B|,D)$ in Section~5 of the SIC manuscript. In every class $B\ne0$. Substitution gives
\[
        f(-\sqrt{13})=D-221-B\sqrt{13}\ne0,
\]
because $D-221$ is rational and $B\sqrt{13}$ is irrational. This excludes the eigenvalue required by the marginal, and hence the proposed local SIC.

\newpage
\section*{The attained minimum in the SIC manuscript}
The original proof of nonexistence~\cite[Sec.~2]{sic} uses algebraic conjugation and purity and requires no spectral enumeration. The same manuscript also determines how closely a state can approach the fifteen required moduli~\cite[Sec.~5]{sic}. With $u=1/\sqrt{13}$ it proves the exact value
\[
  \min_{\psi\in\C^4\otimes\C^3,\ \|\psi\|=1}
      \sum_{P\in\Pset} (|r_P|-u)^2
       =3(1/3-u)^2=\frac{(1-3/\sqrt{13})^2}{3}.
\]
For the lower bound, the quartics give a uniform separation
$|\lambda+\sqrt{13}|\ge\sqrt{13}-3$ for every eigenvalue $\lambda$ of a signed sum. The proof applies this to a kernel vector of $\rho$ and bounds the operator norm of a traceless Hermitian perturbation by its Frobenius norm. Equality is attained by a rank-three marginal with $r_{X\otimes X}=r_{Y\otimes Y}=r_{Z\otimes Z}=1/3$ and $r_P=u$ for each of the other twelve operators. Such a marginal has a purification with a qutrit. The entropy optimization in \emph{StabilizerHalf} and this distance to the local SIC constraints are distinct variational problems.

\paragraph{Earlier uses of the same matrices.}
Liab{\o}tr{\o} studied the same signed Pauli sums in quantum random access codes and reported a numerical search over two-qubit signs, together with the tensor-product spectrum for the all-positive assignment~\cite[Eqs.~(75)--(78) and Table~IV]{liabotro}. Cie\'sli\'nski studied mixed states with all nonidentity Pauli expectations equal in modulus and identified the two-qubit marginals of the three-qubit Hoggar state at modulus $1/3$~\cite[Sec.~V]{cieslinski}. Such a marginal $H$ has rank at most two and hence a purification with a $q=2$ ancilla, while for $q\ge4$ a $q$-dimensional ancilla purifies $wH+(1-w)I/4$ with $w=3/\sqrt{4q+1}$, realizing all fifteen moduli $w/3=1/\sqrt{4q+1}$. The dimension-twelve question fixes the modulus at $1/\sqrt{13}$ and requires a singular marginal.

\paragraph{Chronology and scope of the priority claim.}
The project owner's timestamped archive places completion of the SIC manuscript (\LaTeX{} and PDF) no later than 04:40:32 on 27 November 2026 (Beijing time, UTC+8). The archived source matches the first repository upload byte for byte and contains both the nonexistence proof and the attained minimum above. The recorded Git time of the \emph{StabilizerHalf} commit containing the twelve quartics is 04:42:42 that day, two minutes and ten seconds after this documented completion bound~\cite{half}. The first commit containing the SIC manuscript in its public repository is timestamped 07:52:12~\cite{sic}. To our knowledge, the SIC manuscript gives the first explicit statement and proof of the local dimension-twelve nonexistence result. The cited \emph{StabilizerHalf} submission establishes its two-qubit entropy result and records the quartics. It does not state or derive the local SIC corollary. Its spectral data suffice for the deduction on page~1.

\begingroup\footnotesize
\begin{thebibliography}{9}
\setlength{\itemsep}{2pt}
\bibitem{sic} Chatgpt. \emph{Abelian covariance of symmetric informationally complete measurements in dimension twelve} (2026). \href{https://github.com/rongyulin/physics-open-problems-AI/blob/abc559217923851312de20326ed885d5c5e10f3c/problems/local-sic-dimension-12/LocalSIC12.tex}{Manuscript at initial upload, commit abc5592}.
\bibitem{banacki} M.~Banacki et al. Qubit-qubit-qutrit quantum correlations in $H\to f\bar f V$. \emph{Phys. Rev. D} \textbf{114}, 093007 (2026). \href{https://doi.org/10.1103/hv6p-c34d}{doi:10.1103/hv6p-c34d}.
\bibitem{grassl} M.~Grassl. Tomography of quantum states in small dimensions. \emph{Electron. Notes Discrete Math.} \textbf{20}, 151--164 (2005). \href{https://doi.org/10.1016/j.endm.2005.05.060}{doi:10.1016/j.endm.2005.05.060}.
\bibitem{half} Wangwy-sjtu. \emph{StabilizerHalf}, in \texttt{lean-conjecture-proofs} (2026), commit b39ce7d. \href{https://github.com/Wangwy-sjtu/lean-conjecture-proofs/blob/b39ce7d1828c0682bd2ff7ff9b87c26509bf0eb7/problems/StabilizerHalf/README.md}{Statement and proof outline}. \href{https://github.com/Wangwy-sjtu/lean-conjecture-proofs/blob/b39ce7d1828c0682bd2ff7ff9b87c26509bf0eb7/problems/StabilizerHalf/check.py\#L104-L109}{Exact enumeration and coefficient table}.
\bibitem{liabotro} O.~Liab{\o}tr{\o}. Improved classical and quantum random access codes. \emph{Phys. Rev. A} \textbf{95}, 052315 (2017). \href{https://doi.org/10.1103/PhysRevA.95.052315}{doi:10.1103/PhysRevA.95.052315}.
\bibitem{cieslinski} P.~Cie\'sli\'nski. Pauli flat quantum states mimicking maximal magic (2026). \href{https://arxiv.org/abs/2610.07172}{arXiv:2610.07172}.
\end{thebibliography}
\noindent\copyright\ 2026 Chatgpt. This commentary and its \LaTeX{} source are licensed under \href{https://creativecommons.org/licenses/by/4.0/}{CC BY 4.0}.
\endgroup
\end{document}
